\documentclass[12pt,a4paper]{article}
\usepackage[utf8]{inputenc}
\usepackage[T1]{fontenc}
\usepackage[english]{babel}
\usepackage{amsmath,booktabs,longtable,tabularx,array}
\usepackage{geometry,enumitem,xcolor,fancyhdr}
\usepackage[hyphens]{url}
\usepackage{hyperref}
\geometry{margin=2.5cm}
\hypersetup{colorlinks=true,linkcolor=black,citecolor=blue!60!black,
urlcolor=blue!60!black,
pdftitle={Data Reliability and Representativeness in CNN Channel Gating Research},
pdfauthor={Lidia Cichońska}}
\setlist{itemsep=0.2em,topsep=0.4em}
\renewcommand{\arraystretch}{1.15}
\setlength{\headheight}{15pt}
\setlength{\emergencystretch}{2em}
\pagestyle{fancy}
\fancyhf{}
\lhead{Data Reliability in CNN Research}
\rhead{Lidia Cichońska}
\cfoot{\thepage}
\newcommand{\pp}{pp}
\newcolumntype{Y}{>{\raggedright\arraybackslash}X}

\title{Data Reliability and Representativeness\\
in CNN Channel Gating Research\\[0.5em]
\large Iterative Evaluation of Compression and Adaptive Inference\\
for Components of HCI Systems}
\author{Lidia Cichońska}
\date{6 September 2026}

\begin{document}
\maketitle

\begin{abstract}
This report examines the extent to which the collected experimental data
support conclusions about the accuracy and computational cost of channel
gating in convolutional neural networks. It covers successive experiments
on CIFAR-10, CIFAR-100 and Tiny ImageNet, including data integrity audits,
model diagnostics, statistical comparisons and runtime measurements.
The data support comparisons within a limited benchmark setting and
identify weaknesses in the mechanisms studied. They do not establish
representativeness for a particular user population or readiness for
deployment in a human--computer interaction system. The main limitations
are duplicate images in CIFAR-100, historical monitoring of the test set,
a small number of independent training runs and the absence of data from
the intended application. Negative results motivate subsequent hypotheses
and changes to the evaluation protocol, while the promising early-exit
pilot requires further assessment of accuracy, risk and actual latency.
\end{abstract}

\section{Research Problem and Relevance to HCI}

This report extends the earlier methodological proposal
\emph{Assessing Data Reliability and Completeness in CNN Channel Gating
Methods}.
It retains the six diagnostic markers M1--M6 and the idea of a decision
map for a specialist, adds evidence from subsequent experiments, and
clarifies which analyses have been completed and which remain proposals.
The definitions and interpretations of the markers are updated to reflect
the implemented analyses and their limitations.

The research forms part of doctoral work on deep neural model compression
using structural sparsification, quantization and entropy coding.
This report focuses on learned attention mechanisms and channel gates,
together with their extension towards adaptive network depth.
It does not present a fully validated pipeline combining all three
compression techniques.

The technical question is whether a model can reduce its use of channels
or computational blocks without an unacceptable loss of classification
accuracy. The central methodological question is:
\begin{quote}
Do the provenance, integrity, partitioning and coverage of the data,
together with the experimental design, support a credible assessment
of this trade-off and its generalization to an intended HCI application?
\end{quote}

The intended context is an image recognition component in an interactive
application, for example on a device with limited resources.
Misclassification may provide a person with incorrect information,
while latency may affect the interaction. These are research motivations,
rather than descriptions of an evaluated deployment. The immediate
recipient of the assessment is a researcher or engineer selecting a model
and defining the limits of its use.

In this report, \emph{data reliability} refers to documented provenance
and checks of data correctness; \emph{evaluation rigor} refers to the
ability to reproduce results and characterize their uncertainty;
\emph{completeness} refers to the availability of the observations and
metadata needed to answer the research question; and
\emph{representativeness} refers to the correspondence between the sample
and the population to which a conclusion applies. A dataset may be complete
in terms of files and labels while failing to cover conditions important
to future users.

\section{Data: Provenance, Acquisition and Representativeness}

\subsection{Dataset Characteristics}

The experiments reused public, labeled image datasets. CIFAR-10 and
CIFAR-100 are subsets of the Tiny Images collection \cite{cifar};
Tiny ImageNet is a subset of ImageNet comprising 200 categories
\cite{tinyimagenet}. The project did not collect new images from users
or independently relabel the complete datasets. Inter-annotator agreement
was not assessed. Reference labels were adopted from the benchmarks,
which supports comparison but does not guarantee that every annotation
is semantically correct.

\begin{table}[htbp]
\centering\small
\caption{Image datasets used in the experiments. Counts refer to the official splits.}
\begin{tabularx}{\textwidth}{@{}lrrrY@{}}
\toprule
Dataset & Classes & Training & Evaluation & Format and samples per class \\
\midrule
CIFAR-10 & 10 & 50\,000 & 10\,000 & $32\times32$ RGB;
5\,000 training and 1\,000 test images per class. \\
CIFAR-100 & 100 & 50\,000 & 10\,000 & $32\times32$ RGB;
500 training and 100 test images per class. \\
Tiny ImageNet & 200 & 100\,000 & 10\,000 & $64\times64$ RGB;
500 training and 50 validation images per class. \\
\bottomrule
\end{tabularx}
\end{table}

CIFAR was loaded from existing copies of the official dataset archives;
no additional images were downloaded during the audit.
Tiny ImageNet training images and the official validation annotations
were loaded using a common mapping between class names and label indices.
Validation images were converted to RGB.
The audit documents the state of the local copy but does not provide
its complete download history, so the audit date should not be treated
as the acquisition date. The reported Tiny ImageNet results concern the
labeled validation set or its subsets, rather than the official test set,
whose labels were unavailable in this protocol.

Training used random cropping with padding of 4 pixels for CIFAR or
8 pixels for Tiny ImageNet, together with horizontal flipping.
After conversion to tensors, images were normalized according to the
relevant protocol; Tiny ImageNet used the ImageNet means and standard
deviations. Evaluation did not include random augmentation.
Augmentation increases training variation but does not create new
independent observations or replace data from the target domain.

\subsection{Data Splits and Independence of Evaluation}

The protocols evolved as problems were diagnosed. Safeguards introduced
later should not be attributed retrospectively to earlier experiments.
\begin{enumerate}
\item In the first stage of the research, existing checkpoints were reevaluated on
the complete CIFAR-10 test set. The audit covered seed 42 and 10\,000
images. It was not an independent replication of training.
\item In the 200-epoch global UCLA experiment on CIFAR-100, test performance
was monitored during training. The comparison used the predetermined
epoch 200, which limits favorable checkpoint selection but does not remove
the risk that observed test results influenced subsequent research decisions.
\item A random, deterministic split of the CIFAR-100 training set into
45\,000 training and 5\,000 validation images was introduced later,
using split seed 42. The code permutes indices without guaranteeing
class stratification. Disjoint indices do not exclude identical image content.
\item In the early-exit pilot, the Tiny ImageNet validation set was divided
into the first 5\,000 samples for threshold selection and the remaining
5\,000 for evaluation (\emph{held out from threshold selection}, meaning that these images
were set aside and not used to choose the thresholds).
The adapter sorts samples by filename. This deterministic partition
involves neither randomization nor stratification; class counts and
the distributions of both subsets require inspection.
It is not a new test set from an independent domain.
\end{enumerate}

In future iterations, data partitions and model selection rules should
be fixed before evaluation. Thresholds should be selected on the calibration
set and final accuracy measured on data not used for that selection.
Repeatedly analyzing the same held-out set while developing a method
gradually weakens its independence.

\subsection{Integrity Audit and Its Limitations}

The audit script computed SHA-256 hashes of raw image content before
augmentation and checked sample counts, format, value and label ranges,
and hashes shared across splits. The results are presented in the table below.

\begin{table}[htbp]
\centering\small
\caption{Audit of the local CIFAR copies, conducted on 21 July 2026.}
\begin{tabularx}{\textwidth}{@{}Yrr@{}}
\toprule
Check & CIFAR-10 & CIFAR-100 \\
\midrule
Duplicate image copies beyond the unique count in training & 0 & 14 \\
Duplicate image copies beyond the unique count in testing & 0 & 2 \\
Distinct hashes shared between training and testing & 0 & 10 \\
All classes present and equally frequent within each official split & yes & yes \\
Labels within the valid range & yes & yes \\
\bottomrule
\end{tabularx}
\end{table}

The number of shared hashes need not equal the number of affected records,
because an image may have multiple copies. The CIFAR-100 finding warrants
reporting content leakage and performing a sensitivity analysis that
excludes the affected test records alongside evaluation on the canonical
test set. The procedure for this analysis has been implemented, but
its results have not been established in the evaluation reported here.
The results are therefore not presented as obtained after deduplication.
Future training--validation partitions should keep identical and closely
related images on the same side of the split.

The absence of identical hashes in CIFAR-10 does not establish full
independence: the audit does not detect similar views, resized versions
or photographs from a shared source. A valid label index does not establish
annotation correctness. An equivalent complete integrity audit of
Tiny ImageNet has not been documented.

\subsection{What Do These Data Represent?}

The datasets were selected to enable efficient, controlled comparisons
of methods for classifying small images. Moving from 10 to 100 and 200
classes broadens the task, but does not constitute random sampling of
a future user population. Equal class counts facilitate accuracy comparisons;
they do not reproduce real application frequencies.

The available documentation does not characterize the distribution of
devices, lighting conditions, locations, image quality, rare events or
interaction contexts. Low resolution limits the observation of details
that may determine errors in an application.
In this study, the benchmarks provide no evidence about user tasks,
the consequences of mistakes or the representation of social groups.
Representativeness is therefore conditional: the data are suitable for
evaluation within these benchmarks, but representativeness for a specific
HCI application has not been established.

\section{Experiments and Iterative Method Development}

The main performance measure is Top-1 accuracy: the proportion of images
for which the highest-scoring class matches the reference label.
A checkpoint is a saved state of a trained model, and a seed initializes
the random number generator used in a training run or data split.
Accuracy differences are reported in percentage points (\pp),
which are distinct from relative percentage changes.

\subsection{What Counts as Reduced Computation?}

Gating channel $c$ of a feature map can be written as
\begin{equation}
z'_c=g_c(x)z_c.
\end{equation}
A static mechanism uses the same gate for every image, whereas a dynamic
mechanism conditions the gate on the input. Scaling or zeroing an activation
does not remove a convolution that has already been executed.
Structural sparsification requires physically removing channels and the
corresponding parameters; conditional execution requires skipping operators,
for example later blocks after an earlier classification decision.

Several distinct cost proxies were used in the project.
The proportion of small gates describes candidates for removal,
the budget of 40 out of 64 groups specifies how many groups are selected,
and the early-exit cost reflects the proportion of residual blocks executed.
None of these measures automatically equals the fraction of FLOPs,
runtime or energy consumed. Groups and blocks at different stages
need not have equal computational costs.

\subsection{Channel Gating on CIFAR-10}

The full audit of existing ResNet-18 checkpoints produced the results
below. Differences were calculated using paired predictions on the same
images, with the standard model as the reference.

\begin{table}[htbp]
\centering\small
\caption{CIFAR-10: seed 42, 10\,000 test images; 95\% bootstrap confidence
intervals for the difference, using 1\,000 resamples. Results from the present research.}
\begin{tabularx}{\textwidth}{@{}Yrrr@{}}
\toprule
Variant & Top-1 [\%] & Difference [\pp] & 95\% CI [\pp] \\
\midrule
Standard ResNet-18 & 93.48 & --- & --- \\
Static L1, $\beta=0.2$ & 91.58 & $-1.90$ & $[-2.43,-1.42]$ \\
Static L1, $\beta=0.5$ & 91.28 & $-2.20$ & $[-2.70,-1.73]$ \\
SE, $\beta=0.2$ & 91.49 & $-1.99$ & $[-2.45,-1.49]$ \\
Gumbel, $\beta=0.2$ & 92.60 & $-0.88$ & $[-1.35,-0.43]$ \\
\bottomrule
\end{tabularx}
\end{table}

All gated variants had lower accuracy for these checkpoints.
The intervals exclude zero, but characterize uncertainty associated
with the test sample, rather than variation from retraining.
They do not establish that attention methods necessarily reduce accuracy.

Aggregate accuracy concealed class differences. For Static L1 with
$\beta=0.5$, accuracy for \emph{cat} decreased from 87.1\% to 81.3\%,
a loss of 5.8~\pp, whereas accuracy for \emph{truck} increased from
94.4\% to 95.0\%. Both classes had 1\,000 test samples.
This demonstrates an uneven change in benchmark accuracy; without further
causal analysis, it does not establish class underrepresentation or
harm to a particular group of people.

In the audit of these same models, the dense convolution operation count
did not decrease for the static variants.
Increasing the proportion of small gates therefore did not demonstrate
an actual reduction in execution cost. This motivated separating gate
diagnostics from claims about compression benefits in deployment.

\subsection{From Uncertainty to the Value of Additional Computation}

In the next stage of the research, the hypothesis became more specific:
with exactly 40 out of 64 ResNet-50 channel groups retained, uncertainty
was expected to improve group allocation across stages.
After 200 epochs on CIFAR-100 with seed 42, global UCLA achieved
69.82\% Top-1 accuracy, compared with 71.58\% for the utility-only control.
The paired difference was $-1.76$~\pp, with a 95\% CI of
$[-2.58,-0.93]$~\pp.
The control retained the same backbone, group count and protocol;
this does not guarantee identical actual runtime.

The correlation between predicted uncertainty and classification error
was 0.054. Diagnostics indicated that the supervision signal derived
from training errors diminished as the model fitted the training data.
Moreover, a single classification correctness label did not specify
which stage would benefit from additional computation.
This particular implementation was judged unsuccessful.
The conclusion is limited by the single seed, test monitoring and
the duplicate-image findings.

The next hypothesis concerned \emph{value-of-compute}: the model was to
predict the value of computation at a stage from the counterfactual
change in loss after restricting that stage's budget.
A training--validation--test split was introduced for CIFAR-100 and
the allocator was improved. On Tiny ImageNet, an initial signal after
20 epochs did not persist in the longer, 100-epoch protocol:
\begin{center}
\begin{tabular}{@{}rrrr@{}}
\toprule
Seed & Value-of-compute [\%] & Utility-only [\%] & Difference [\pp] \\
\midrule
42 & 58.36 & 58.33 & $+0.03$ \\
123 & 58.74 & 59.33 & $-0.59$ \\
Mean & 58.55 & 58.83 & $-0.28$ \\
\bottomrule
\end{tabular}
\end{center}
These are validation results. The utility-only control for seed 42 uses
a completed rerun; the first run was documented only up to epoch 23.
Two seeds do not support a precise estimate
of the effect distribution, but provide no basis for claiming superiority.
The change in the conclusion between 20 and 100 epochs illustrates why
a promising pilot cannot replace evaluation after longer training.

An \emph{oracle} audit, with access to reference labels and legal group
transfers, found a potential Top-1 improvement of approximately 11~\pp\
on 512 images. This diagnoses the decision space with access to the answer;
it is not a result achievable under ordinary inference conditions.
Policies predicting actions from features after the initial convolutional
module and the first and second residual stages did not exceed the accuracy
of the majority-action baseline on held-out data.
The small sample and restricted transfer space further limit interpretation.

\subsection{Negative Result for Semantic Feedback}

A further experiment tested whether deep features used only during training
could improve learned static channel selection.
Both variants used Tiny ImageNet, ResNet-50, 10 out of 16 first-stage groups,
AdamW and 20 epochs with seed 42.
Final-epoch accuracy on the deployment path was 49.83\% with semantic
feedback and 55.19\% for the control, a difference of $-5.36$~\pp.
Feedback was disabled during evaluation.

A plausible explanation is the mismatch between image-dependent group
selection during training and a single ranking at deployment.
This is a hypothesis about the failure mechanism, not an independently
established causal relationship. Further scaling of this version was
discontinued; one seed does not justify rejecting all feedback methods.
The control result also does not demonstrate a speedup, since it uses
activation masking.

\subsection{Adaptive Depth: A Promising but Incomplete Pilot}

Following the channel-routing results, a further hypothesis was tested:
easier images could complete classification after the second or third
stage of ResNet-50. Three classifier heads were trained jointly for
20 epochs on Tiny ImageNet.
Confidence thresholds were selected on the calibration subset to maximize
accuracy subject to a target mean cost of no more than 85\%.
The exit costs were $7/16$, $13/16$ and $1$, respectively,
based on the number of residual blocks executed.

\begin{table}[htbp]
\centering\small
\caption{Early exit on 5\,000 images held out from threshold selection.
Full depth denotes the final exit of the same jointly trained model.
SD was calculated across three seeds using the sample variance denominator
$n-1$; it is not a confidence interval.}
\begin{tabularx}{\textwidth}{@{}lYYYY@{}}
\toprule
Seed & Early exit [\%] & Full depth [\%] & Difference [\pp] & Block cost [\%] \\
\midrule
42 & 54.96 & 54.96 & $0.00$ & 84.75 \\
123 & 53.46 & 53.78 & $-0.32$ & 85.08 \\
2026 & 54.12 & 54.02 & $+0.10$ & 84.95 \\
Mean & 54.18 & 54.25 & $-0.07$ & 84.92 \\
SD & 0.75 & 0.62 & 0.22 & 0.17 \\
\bottomrule
\end{tabularx}
\end{table}

The mean result is consistent with a useful trade-off, but does not
establish statistical equivalence in accuracy.
It is also not a comparison with an independently trained full-depth
model without auxiliary heads.
For seed 123, held-out cost exceeded the target by 0.08~\pp;
a calibration constraint does not guarantee the same cost on new data.

A control trained to always terminate at stage three achieved 53.55\%
for seed 42 and 53.67\% for seed 123, with a block cost of 81.25\%.
However, its evaluation script uses the full validation set of 10\,000
images, whereas the adaptive policy above is evaluated on 5\,000.
The costs also differ. This juxtaposition is descriptive and does not
establish superiority of either method.
Evaluation on identical samples, comparison of accuracy--cost curves
and a third replication of the static control are required;
a completed result for the latter was not found.

\subsection{What Did the Runtime Measurement Establish?}

For seed 42, a CUDA benchmark measured execution of the network prefix
ending at a specified stage, using batch size 1, 20 warm-up iterations
and 60 repetitions. Each measured forward pass was followed by CUDA
synchronization. The script uses a synthetic image tensor and specifies
the exit in advance.
Median times were 2.63~ms, 4.62~ms and 5.65~ms, with corresponding
P95 values of 2.94~ms, 5.06~ms and 5.79~ms.

Held-out exit frequencies were 15.06\%, 36.18\% and 48.76\%.
The weighted sum of prefix medians gives the approximation
\begin{equation}
\widetilde{T}=0.1506\cdot2.63+0.3618\cdot4.62+
0.4876\cdot5.65\approx4.82\ \mathrm{ms}.
\end{equation}
This is approximately 14.6\% below the median for the full prefix.
However, \emph{a weighted sum of medians is not the median of the policy's
runtime distribution}, nor is it a direct measurement of its mean.
The benchmark confirms shorter execution of prefixes that omit later
blocks. It does not measure the complete adaptive decision process,
input preprocessing, data transfer, queueing or the interface.
It therefore does not support a claim that the entire application
is 14.6\% faster.
Reproduction also requires the exact GPU model and software versions;
stating only that CUDA was used is insufficient.

\section{Verification of Correctness and Evaluation Rigor}

\subsection{Data Auditing and Diagnostic Markers M1--M6}

Six markers of model behavior were developed in the first stage of the research.
The integrity audit complements them; it is not one of the original
M1--M6 markers. These markers are not six independently validated measures
of data representativeness. Their scope is summarized below.

\begingroup\small
\begin{longtable}{@{}p{0.23\textwidth}p{0.34\textwidth}p{0.35\textwidth}@{}}
\caption{Diagnostic markers: measurements, evidence and limitations.}\\
\toprule
Marker & Measurement and available evidence & Limits of interpretation \\
\midrule
\endfirsthead
\toprule
Marker & Measurement and available evidence & Limits of interpretation \\
\midrule
\endhead
M1: per-class accuracy & Accuracy, sample counts and changes relative
to a control for each class; audited on CIFAR-10. & Class results require
uncertainty estimates and replication; they do not describe every
context within a class. \\
M2: error concentration & The distribution of errors relative to class
frequencies; evaluated on CIFAR-10. & Concentration may reflect class
difficulty. It does not establish that compression is the cause
or that discrimination has occurred. \\
M3: threshold stability & Gate distributions and the proportion of positive
values below $10^{-2}$; approximately 99.9--100\% for two L1 variants
in the marker audit. & Describes gate values, rather than physically
removed channels, variation across training seeds or image
representativeness. \\
M4: gate--class dependence & Analysis of input-dependent gates;
marked as not applicable for static attention. & Equal class counts
prevent meaningful interpretation of correlation with class frequency.
Mean gate values do not replace measured cost. \\
M5: representation change & Embedding comparisons using cosine similarity
of mean vectors, KL divergence and MMD. Cosine values for two L1 variants:
0.9612 and 0.9558. & Similarity of means does not establish coverage
of the data space or robustness to a domain shift. \\
M6: quantization stability & Simulated quantization of FP32 activations
to 256 levels; reconstruction errors and distribution statistics.
& Does not evaluate the accuracy, memory use or runtime of an actual
INT8 model after PTQ/QAT. \\
\bottomrule
\end{longtable}
\endgroup

The six-marker audit concerns a separate series of trained models
from the CIFAR-10 variant comparison presented earlier.
Reference-model accuracy is 88.51\% in the marker audit and 93.48\%
in the variant comparison. Their measurements must not be combined
into a single result for the same model.
The threshold sensitivity indicator in the marker audit is defined as
\begin{equation}
\mathrm{fragile\_ratio}=
\frac{\#\{g:0<g<10^{-2}\}}{\#\{g:g>0\}}.
\end{equation}
A value close to one concerns the positive gates included in the
denominator; it is not the proportion of all channels physically removed.

Automatic warning labels assigned by the markers are working heuristics.
The absence of a warning is not a certificate of correctness, and missing
data must be distinguished from a measure being inapplicable.
The markers should not be summed into a single score that could conceal
data leakage or domain mismatch.

An assessment card should bring together the intended use, study population,
data sources, status of each audit, results with uncertainty estimates,
and missing controls. This makes it possible to identify which conclusions
are supported, which require a warning, and which cannot yet be drawn.
In this report, the documented results and explicit limitations serve
this purpose; the usefulness of the card to its intended recipients
has not been evaluated in a separate study.

\subsection{Statistics: Unit of Analysis and Scope of Uncertainty}

When models are evaluated on the same images, the unit of analysis is
a pair of predictions. For candidate model $A$ and control $B$,
\begin{equation}
d_i=\mathbf{1}[\hat y_i^A=y_i]-\mathbf{1}[\hat y_i^B=y_i],
\qquad \widehat{\Delta}=\frac{1}{n}\sum_{i=1}^{n}d_i.
\end{equation}
The bootstrap samples shared image indices with replacement and
recalculates $\widehat{\Delta}$.
The CIFAR audits used 1\,000 resamples to obtain 95\% intervals.
Recording discordant prediction pairs helps establish whether one model
corrects the other's errors; separate accuracy intervals do not replace
an interval for their difference.

This bootstrap is conditional on the trained models and assumes
appropriate independence of observations.
It does not account for training randomness, training-data selection
or hyperparameter search.
For related images, resampling groups defined by a common source should
be considered. For repeated training runs, paired differences should
also be reported separately for each seed.
The three early-exit seeds provide an initial estimate of variability;
15\,000 predictions on the same 5\,000 images do not constitute
15\,000 independent images.

Class results require caution: CIFAR-100 has 100 test images per class,
while the full Tiny ImageNet validation set has 50.
Partitioning validation reduces these counts, which need not remain equal.
Testing multiple classes and variants increases the risk of chance
findings; the current class analyses should be treated as diagnostic.
A future study should prespecify its primary comparison and account
for multiplicity in secondary analyses, for example using Holm's correction.

If the maximum acceptable accuracy loss is $\delta$, establishing
noninferiority requires the appropriate lower confidence bound
for the difference to exceed $-\delta$.
Equivalence requires an interval contained within prespecified equivalence
margins. A small mean difference or a nonsignificant result alone
does not meet these conditions.
The 0.5~\pp\ threshold used in the pilot is a criterion for continuing
research, not a validated acceptance margin for HCI users.

\subsection{Implementation Correctness and Reproducibility}

Tests were developed to check data partitions, the budget controller,
legality of oracle actions and early-exit execution.
Their role is to detect technical errors, such as an incorrect group
count or execution of blocks beyond the selected exit.
Passing such tests does not establish the research hypothesis or
the correctness of dataset labels.
This report draws on completed experimental results and analysis of their implementation;
preparing it does not constitute retraining or independent reproduction
of the complete experimental campaign.

Full reproduction of a result requires sample identifiers, data and
checkpoint hashes, the code version, configuration, training and split
seeds, the epoch selection rule, predictions and an environment description.
The research documentation includes experimental settings, training histories,
selected audits and saved model states,
but a complete record of predictions, embeddings and per-image timing
should not be assumed for every experiment.
Missing fields must be collected before the relevant analysis,
rather than replaced with zero.
Results for the best validation checkpoint and the final epoch must
be reported separately.

\section{Methodological Context in the Literature}

Squeeze-and-Excitation demonstrates input-dependent channel weighting
\cite{senet}.
BilevelPruning combines dynamic and static channel pruning, including
permanent removal of some parameters \cite{bilevel},
while Pick-or-Mix studies dynamic channel sampling and its efficiency
\cite{pix}.
These studies motivate separating the role of attention from actual
structural and computational reductions.
Their numerical results are not compared directly with the local pilots
because the protocols and execution conditions differ.

Meronen and colleagues examine overconfidence in dynamic networks
and the importance of modeling uncertainty when deciding on
a computational budget \cite{overconfidence}.
This motivates evaluating calibration and error risk when computation
terminates early. A high softmax value alone is insufficient justification
for trusting a prediction.

Documenting data provenance, composition and intended uses follows
the \emph{Datasheets for Datasets} approach \cite{datasheets}.
Reporting performance across subgroups and defining a model's scope
of use is consistent with \emph{Model Cards} \cite{modelcards}.
The proposed assessment card combines these principles with CNN
compression diagnostics.
Its value for engineering decisions still requires separate validation;
it is not presented as a new, already validated measure of representativeness.
The cited studies provide selected methodological reference points,
rather than an exhaustive literature review.

\section{Subsequent Hypotheses and Verification Plan}

The results changed both the modeling direction and the evaluation
methodology.
Error uncertainty was followed by the question of the value of additional
computation; oracle results were examined in terms of predictability
without access to labels; and channel masking was supplemented by
an experiment on executable block skipping.
The next iterations should address specific gaps in the evidence:

\begin{enumerate}
\item \emph{Integrity hypothesis:} the ranking of models on CIFAR-100
will remain unchanged after excluding test records that overlap
with training data.
The analysis should report affected indices, class counts and paired
changes in performance, while retaining evaluation on the canonical test.
This is a planned sensitivity analysis, not a claim that its outcome
is already known.
\item \emph{Adaptation hypothesis:} early exit will offer a useful
trade-off relative to fixed-depth models at comparable measured latency.
The evaluation samples should be aligned, the seed-2026 control completed,
and uncertainty in paired differences estimated.
Longer training is justified after this control.
\item \emph{Coverage hypothesis:} images poorly represented in the
training data will more frequently produce additional errors under
the adaptive policy.
A distance measure from training examples in embedding space can be
prespecified, and accuracy, calibration and exit decisions compared
across regions of different density.
Regions should be defined using training or validation data and
the hypothesis tested on held-out observations.
\item \emph{Robustness hypothesis:} the selected policy will maintain
acceptable risk under blur, noise and changes in image quality.
Synthetic perturbation tests should follow a fixed protocol and
be supplemented with data from the target device.
PCA/UMAP projections may assist interpretation, but do not replace
quantitative evaluation or domain-specific data.
\item \emph{HCI usefulness hypothesis:} reducing computation will
improve human task performance without increasing harmful errors.
The complete execution path should first be measured on the target
device, including P50, P95 and P99 latency, followed by a user study
with a defined population, task and acceptance criteria.
\end{enumerate}

The sample size for the planned noninferiority study should be selected
using the expected number of discordant prediction pairs and the maximum
acceptable accuracy loss.
In a user study, the independent unit may be the participant rather
than each click or task repetition.
Sample size and analysis should reflect that structure.
These analyses are not claimed as part of the experiments completed
to date.

\section{Social Responsibility and Limits of Use}

Compression may make image recognition more accessible on less capable
hardware, but its accuracy costs may affect some usage conditions
more than others.
For example, poor lighting or a lower-quality camera may change
recognition difficulty.
The current data do not establish who would be affected by such
limitations; a particular social distribution of effects should not
be assumed without evidence.

Image classes do not constitute validated social groups.
Per-class analysis diagnoses performance but is not a complete fairness
audit. An attention weight alone is also not a human-understandable
causal explanation of a decision.
Recipients of the assessment should be informed about test conditions,
uncertainty, possible consequences of errors and uses excluded
from the model's evaluated scope.

Before an application involving people is considered, data and criteria
must be defined with a domain expert and appropriately selected
participants.
A plan for collecting new data should address informed consent where
applicable, data minimization, protection of identifiable images
and access rules.
Using a public benchmark does not replace checking its terms of use
or assessing whether its data are suitable for transfer to a new domain.

For uncertain inputs, possible responses include running the full
computational path, presenting a warning or referring the case to
a person. Each option requires evaluation: the full network can also
make mistakes, and a warning may not be understood.
No study has yet assessed the effects of these mechanisms on task
completion time, cognitive workload or the appropriateness of user trust.
Energy has not been measured either, so a lower cost proxy is not
interpreted as evidence of energy savings.

\section{Answer to the Reliability and Representativeness Question}

The collected data provide a documented basis for comparing particular
checkpoints and diagnosing problems in the studied methods on small
image benchmarks.
They support the conclusion that the evaluated dynamic gating variants
did not demonstrate superiority over the corresponding controls,
and that small gate values do not establish execution savings.
The replicated early-exit signal justifies further evaluation,
but does not yet establish equivalence in accuracy, superiority
over fixed depth or acceleration of the complete application.

The credibility of these conclusions has explicit limits:
CIFAR-100 duplicates, historical test monitoring, few training seeds,
different evaluation subsets in some comparisons and incomplete
measurement of deployment conditions.
Representativeness for a specific HCI system remains unestablished,
because the target population of images, devices, tasks and users
has not been defined and studied.

The contribution of the iterations completed so far is both to narrow
the technical hypotheses and to improve the evaluation process.
The next stage should supply the missing data and controls so that
decisions about model use follow from meaningful evidence and
an understanding of its limitations.

\begin{thebibliography}{99}
\bibitem{cifar} A. Krizhevsky, \emph{Learning Multiple Layers of Features
from Tiny Images}, 2009. Dataset documentation:
\url{https://cave.cs.toronto.edu/kriz/cifar.html}.
\bibitem{tinyimagenet} Y. Le, X. Yang, \emph{Tiny ImageNet Visual Recognition
Challenge}, Stanford CS231n, 2015.
\url{https://cs231n.stanford.edu/reports/2015/pdfs/yle_project.pdf}.
\bibitem{senet} J. Hu, L. Shen, G. Sun,
\emph{Squeeze-and-Excitation Networks}, CVPR, 2018, pp. 7132--7141.
\url{https://openaccess.thecvf.com/content_cvpr_2018/html/Hu_Squeeze-and-Excitation_Networks_CVPR_2018_paper.html}.
\bibitem{bilevel} S. Gao, Y. Zhang, F. Huang, H. Huang,
\emph{BilevelPruning: Unified Dynamic and Static Channel Pruning for
Convolutional Neural Networks}, CVPR, 2024, pp. 16090--16100.
\url{https://openaccess.thecvf.com/content/CVPR2024/html/Gao_BilevelPruning_Unified_Dynamic_and_Static_Channel_Pruning_for_Convolutional_Neural_CVPR_2024_paper.html}.
\bibitem{pix} A. Kumar, D. Kim, J. Park, L. Behera,
\emph{Pick-or-Mix: Dynamic Channel Sampling for ConvNets}, CVPR, 2024,
pp. 5873--5882. \url{https://arxiv.org/abs/2406.10935}.
\bibitem{overconfidence} L. Meronen et al.,
\emph{Fixing Overconfidence in Dynamic Neural Networks}, WACV, 2024.
\url{https://arxiv.org/abs/2302.06359}.
\bibitem{datasheets} T. Gebru et al., \emph{Datasheets for Datasets},
Communications of the ACM, 64(12), 2021, pp. 86--92.
\url{https://arxiv.org/abs/1803.09010}.
\bibitem{modelcards} M. Mitchell et al., \emph{Model Cards for Model Reporting},
FAT*, 2019, pp. 220--229. \url{https://arxiv.org/abs/1810.03993}.
\end{thebibliography}
\end{document}
